% Part: computability
% Chapter: recursive-functions
% Section: primitive-recursion

\documentclass[../../../include/open-logic-section]{subfiles}

\begin{document}

\olfileid{cmp}{rec}{pre}
\olsection{بنیادی بازگشت}

قدرتی عدداں دی اک خاصیت ایہ اے کہ ہر قدرتی عدد تک~$0$
توں جانشین بنانے دا عمل~$+1$ محدود واری لا کے پہنچیا جا
سکدا اے---ہر قدرتی عدد یا تاں~$0$ اے یا جانشین اُتے
جانشین \dots{} لا کے~$0$ توں بنیا اے۔ ایس خاصیت
نوں ورت کے فنکشن
~$h\colon\Nat \to \Nat$ متعین کرن دا اک طریقہ ایہ اے:
(a)~دسّو کہ $h$ دی دلیل~$0$ لئی قیمت کیہ اے، تے
(b)~ایہ وی دسّو کہ $h(x)$ دی قیمت معلوم ہووے تاں
$h(x+1)$ دی قیمت کِویں حساب کرنی اے۔ (a) ساڈے نوں
سِدھا $h(0)$ دَسدا اے، ایس لئی $h$~،~$0$ اُتے
متعین اے۔ ہن (b) دی ہدایت $x=0$ لئی لا کے اسی
$h(1) = h(0+1)$ نوں~$h(0)$ توں حساب کر سکدے آں۔
اوہی ہدایت $x=1$ لئی ورت کے اسی $h(2) = h(1+1)$
نوں~$h(1)$ توں حساب کردے آں، تے ایویں اگے ودھدے آں۔
ہر قدرتی عدد~$x$ لئی آخر اوہ قدم آ جاندا اے جتھے
$h(x)$ توں $h(x+1)$ دی قیمت متعین کردے آں؛ ایس
لئی $h(x)$ ہر $x \in \Nat$ لئی متعین اے۔

مثال لئی، فرض کرو اسی $h\colon \Nat \to \Nat$ نوں
ہیٹھ لکھیاں دو مساواتاں نال متعین کردے آں:
\begin{align*}
h(0) & =  1\\
h(x+1) & =  2 \cdot h(x)
\end{align*}
جے ضرب کرنی پہلے توں آؤندی ہووے، تاں ایہ مساواتاں
اوپرلے (a) تے~(b) لئی لوڑ دی معلومات دیندیاں نیں۔
دوجی مساوات نوں وار وار لا کے سانوں ملدا اے:
\begin{align*}
  h(1) & = 2\cdot h(0) = 2,\\
  h(2) & = 2\cdot h(1) = 2\cdot 2,\\
  h(3) & = 2 \cdot h(2) = 2\cdot 2 \cdot 2,\\
  & \vdots
\end{align*}
اسی ویکھدے آں کہ ایہ متعین کیتا فنکشن~$h$،
$h(x) = 2^x$ اے۔

قدرتی عدداں دی ایس خاصیت نال یقین ہو جاندا اے کہ
ایہ دو شرطاں پوری کرن والا صرف اک ہی فنکشن~$h$ اے۔
ایہو جیہیاں دو مساواتاں نوں فنکشن~$h$ دی
\emph{بنیادی بازگشت دے ذریعے تعریف (definition by primitive
  recursion)} آکھدے نیں۔ اِس نوں «بازگشتی» ایس لئی
آکھدے نیں کہ اسی~$h$ نوں ایس طرح تعریف کردے آں
کہ تعریف، خاص طور اُتے دوجی مساوات، دے سجے پاسے
$h$ آپ وی آؤندا اے۔ اِس نوں «بنیادی» ایس لئی
آکھدے نیں کہ $h(x+1)$ نوں متعین کرن لئی اسی
صرف~$h(x)$، یعنی اوہدے عین پہلے والی قیمت، ورتدے
آں۔~$\Nat$ اُتے فنکشن نوں بازگشتی طریقے نال تعریف
کرن دا ایہ سب توں سادہ طریقہ اے۔

اسی جمع تے ضرب ورگے ہور وی بنیادی فنکشن بنیادی
بازگشت نال تعریف کر سکدے آں۔ پر انہاں صورتاں وچ
فنکشن $2$-جگہ والے نیں۔ اسی دلیل دی اک جگہ مقرر
کر لَیندے آں تے دوجی اُتے بازگشت چلاوندے آں۔ مثال
لئی $\Add(x, y)$ دی تعریف لئی اسی~$x$ مقرر کرکے
پہلاں $y=0$ لئی قیمت، تے فیر $y+1$ لئی قیمت
~$y$ دے حوالے نال متعین کردے آں۔ چونکہ $x$ مقرر
اے، ایس لئی اوہ تعریف کرن والیاں مساواتاں دے کھبے
تے سجے دوہاں پاسیاں آوے گا:
\begin{align*}
\Add(x,0) & =  x\\
\Add(x,y+1) & =  \Add(x,y)+1
\end{align*}
ایہ مساواتاں $\Add$ دی قیمت سارے $x$
\emph{تے}~$y$ لئی متعین کردیاں نیں۔ مثال لئی
$\Add(2,3)$ کڈھن واسطے اسی $x = 2$ رکھ کے
تعریفی مساواتاں ورتدے آں: پہلی توں
$\Add(2,0) = 2$ لبھدا اے، فیر دوجی وار وار
لا کے $\Add(2,1) = 2 + 1 = 3$،
$\Add(2, 2) = 3 + 1 = 4$،
$\Add(2, 3) = 4 + 1 = 5$ لبھدا اے۔

$\Add$ دی تعریف وچ اسی دوجی مساوات دے سجے
پاسے $+$ ورتیا، پر صرف~$1$ جوڑن لئی۔ دوجے لفظاں
وچ اسی جانشین فنکشن $\Succ(z) = z+1$ نوں پچھلی
قیمت $\Add(x,y)$ اُتے لا کے $\Add(x,y+1)$
متعین کیتا۔ ایس لئی بازگشتی تعریف نوں اکّو فنکشن
دے ذریعے سمجھ سکدے آں، جیہڑا اسیں پچھلی قیمت اُتے
لاؤندے آں۔ پر فنکشن نوں صرف پچھلی قیمت اُتے ہی
نہیں، بلکہ $x$ تے~$y$ اُتے وی منحصر ہون دینا
نقصان دہ نہیں---تے کدے ضروری ہوندا اے۔ ویکھو:
\begin{align*}
  \Mult(x,0) & =  0 \\
  \Mult(x,y+1) & =  \Add(\Mult(x,y),x)
\end{align*}
ایہ فنکشن $\Mult$ دی بنیادی بازگشتی تعریف اے:
فنکشن $\Add$ نوں پچھلی قیمت $\Mult(x,y)$
تے پہلی دلیل~$x$ دوہاں اُتے لایا جاندا اے۔
ایہ فنکشن~$\Mult(x,y)$ نوں سارے $x$ تے~$y$
لئی وی متعین کردا اے۔ مثال لئی $\Mult(2,3)$،
$\Mult(2,0)$، $\Mult(2,1)$،
$\Mult(2,2)$ تے~$\Mult(2,3)$ وار وار حساب
کرن نال متعین ہوندا اے:
\begin{align*}
  \Mult(2,0) & = 0\\
  \Mult(2,1) & = \Mult(2,0+1) =
  \Add(\Mult(2,0), 2) = \Add(0, 2) = 2\\
  \Mult(2,2) & = \Mult(2,1+1) =
  \Add(\Mult(2,1), 2) = \Add(2, 2) = 4\\
  \Mult(2,3) & = \Mult(2,2+1) =
  \Add(\Mult(2,2), 2) = \Add(4, 2) = 6
\end{align*}

تاں عام نمونہ ایہ اے: فنکشن~$h(x_0, \dots, x_{k-1}, y)$
دی بنیادی بازگشتی تعریف لئی اسی دو مساواتاں دیندے
آں۔ پہلی $h(x_0, \dots, x_{k-1}, 0)$ دی قیمت
~$h$ دا حوالہ دتے بغیر متعین کردی اے۔ دوجی،
$h(x_0, \dots, x_{k-1}, y+1)$
دی قیمت $h(x_0, \dots, x_{k-1}, y)$،
باقی دلیلاں $x_0$، \dots،~$x_{k-1}$ تے~$y$
دے حوالے نال متعین کردی اے۔ ایس دوجی مساوات وچ
~$h$ دی صرف عین پہلے والی قیمت ورتی جا سکدی اے۔
جے اسی انہاں دو مساواتاں دے سجے پاسے دتے عملاں
نوں خود فنکشن $f$ تے~$g$ منّیے، تاں بنیادی بازگشت
نال نواں فنکشن~$h$ تعریف کرن دا عام نمونہ ایہ اے:
\begin{align*}
  h(x_0, \dots, x_{k-1}, 0) & = f(x_0, \dots, x_{k-1})\\
  h(x_0, \dots, x_{k-1}, y+1) & =
  g(x_0, \dots, x_{k-1}, y, h(x_0, \dots, x_{k-1}, y))
\end{align*}
$\Add$ دی صورت وچ $k=1$ تے $f(x_0) = x_0$
(شناختی فنکشن)، تے $g(x_0, y, z) = z + 1$
(اوہ $3$-جگہ والا فنکشن جیہڑا اپنی تیسری دلیل
دا جانشین واپس کردا اے):
\begin{align*}
  \Add(x_0, 0) & = f(x_0) = x_0\\
  \Add(x_0, y+1) & = g(x_0, y, \Add(x_0, y)) =
  \Succ(\Add(x_0, y))
\end{align*}
$\Mult$ دی صورت وچ $f(x_0) = 0$ (اوہ مستقل
فنکشن جیہڑا ہمیشہ~$0$ واپس کردا اے) تے
$g(x_0, y, z) = \Add(z,x_0)$ (اوہ
$3$-جگہ والا فنکشن جیہڑا اپنی آخری تے پہلی
دلیل دا جوڑ واپس کردا اے):
\begin{align*}
  \Mult(x_0,0) & =  f(x_0) = 0 \\
  \Mult(x_0,y+1) & =  g(x_0, y, \Mult(x_0,y)) =
  \Add(\Mult(x_0,y), x_0)
\end{align*}

\end{document}
